\begin{abstract}
HTTP security classification is an attractive use case for compact language models because requests contain a mixture of structured fields, free-form payloads, authorization context, and increasingly, natural-language instructions that may themselves be adversarial. However, reported gains can be misleading when test labels leak into prompt design, when structured-output compliance is conflated with security judgment, or when evaluation uses only synthetic in-distribution samples. We present \model{}, a 4B-parameter language model adapted with QLoRA for request-side HTTP security classification into \texttt{safe}, \texttt{attack}, or \texttt{unknown}, together with a leakage-controlled evaluation protocol. The primary evaluation uses a frozen 413-sample authorized LAN challenge set that was excluded from training, prompt selection, threshold selection, and model selection. Reference labels were produced through dual-AI annotation and independent third-model adjudication, audited against model-visible context, and subsequently confirmed by a confidential blind human review covering all 413 samples with zero label changes. We compare eight systems: a frozen repository rule engine, a character n-gram logistic-regression baseline, Qwen3-4B zero-shot and fixed few-shot baselines, and four \model{} inference artifacts. The HF-LoRA model obtains a verdict Macro-F1 of 0.4658 versus 0.4175 for Qwen3-4B zero-shot, a paired-bootstrap difference of +0.0482 with 95\% CI [+0.0113,+0.0835]. Attack recall rises from 0.7113 to 0.9072, while safe false-positive rate falls from 0.3206 to 0.2057. These gains are nevertheless insufficient for autonomous enforcement: the model misclassifies 43 of 209 safe requests as attacks and never emits \texttt{unknown} for any of 107 reference-unknown samples. A Q4\_K\_M artifact reduces measured storage from approximately 7.49 to 2.33~GiB while retaining a Macro-F1 of 0.4559. Controlled inference-path diagnostics further show that changing message structure alters 11 of 413 verdicts, whereas byte-identical serialized prompts yield identical verdicts across the tested Transformers merged-BF16 and Ollama-f16 paths. The results support compact-model task adaptation for attack-sensitive shadow analysis, but also expose abstention failure, distribution shift, and deployment risks that aggregate accuracy alone would conceal.
\end{abstract}

\begin{IEEEkeywords}
HTTP security, intrusion detection, large language models, QLoRA, prompt injection, abstention, leakage-controlled evaluation, quantization.
\end{IEEEkeywords}

\section{Introduction}
\label{sec:introduction}

Application-layer security systems increasingly process evidence that does not fit cleanly into a fixed signature. An HTTP request can combine a conventional injection payload, an authorization-sensitive object identifier, ambiguous business context, and natural-language content intended for an AI agent. Signature engines remain fast and auditable, but their coverage depends on rules written before deployment. Statistical classifiers can generalize beyond exact signatures, yet they often require stable feature distributions and do not naturally express structured decisions or uncertainty. Large language models (LLMs) offer a third design point: they can jointly interpret paths, headers, bodies, and textual context, and can emit a schema suitable for downstream policy engines. Their instruction-following behavior, however, also creates a new attack surface because untrusted request content may attempt to redirect the classifier itself~\cite{perez2022promptinject,greshake2023indirect,schulhoff2023hackaprompt}.

This paper studies a deliberately narrow question: \emph{can a compact, locally deployable language model improve request-side HTTP security classification over its untuned base model and conventional baselines under a frozen, leakage-controlled external challenge protocol?} We focus on a 4B-parameter model rather than a frontier-scale service because local deployment can reduce data exposure, model-call cost, and operational dependency. The base model is Qwen3-4B-Instruct-2507~\cite{qwen3,qwen2507}; task adaptation uses low-rank adapters and 4-bit quantized training following LoRA and QLoRA~\cite{hu2021lora,dettmers2023qlora}.

The evaluation problem is more difficult than reporting a single accuracy number. First, security datasets are vulnerable to direct and indirect leakage: labels may influence prompt wording, examples, thresholds, rules, or checkpoint selection. Second, an LLM may appear weak merely because it violates a strict JSON schema even when its semantic judgment is useful; conversely, perfect schema compliance does not imply correct security classification. Third, real request-side evidence is sometimes insufficient. A secure system should be able to abstain rather than force every request into a binary allow/block interpretation. Fourth, deployment artifacts and message serialization can change model behavior even when nominal weights and prompts appear equivalent.

To address these concerns, we separate prediction from scoring, freeze the protocol before model execution, preserve raw outputs, and evaluate \texttt{safe}, \texttt{attack}, and \texttt{unknown} as distinct classes. The 413-sample \challenge{} challenge set is external to all training and few-shot examples. The inference process reads only request-side model inputs, while labels are loaded by a separate scoring process after predictions have been written. Two AI annotators and an independent third AI adjudicator produced the initial references; a model-visible-context audit corrected provenance errors before evaluation. After the frozen evaluation, a confidential blind human review independently confirmed all 413 labels without modifying the label file. Because the review occurred after predictions and produced zero label changes, it strengthens label provenance without introducing test-time tuning.

Our main finding is mixed. Task adaptation materially improves attack-sensitive classification relative to the base model: \model{} increases verdict Macro-F1 and attack F1, improves attack recall to 0.9072, and reduces safe false positives relative to Qwen3-4B zero-shot. Yet the absolute safe false-positive rate remains 20.57\%, and the model's unknown recall is 0. These two failures prevent a production-safe autonomous-blocking interpretation. The traditional rule engine displays the opposite profile: only 1.44\% safe FPR but 11.34\% attack recall. The evidence therefore favors complementary deployment---rules for high-precision known patterns and the language model for shadow-mode recall and analyst prioritization---rather than replacement.

This work makes four contributions:
\begin{enumerate}
    \item \textbf{A compact structured security classifier.} We formulate request-side HTTP analysis as a three-class task with attack-family, action, context-sufficiency, confidence, and schema-validity outputs, and adapt a 4B instruction model using QLoRA.
    \item \textbf{A leakage-controlled evidence protocol.} We provide protocol commits, content hashes, isolated inference and scoring paths, frozen few-shot examples, per-sample predictions, paired statistics, and claim-to-evidence mappings for a 413-sample challenge set.
    \item \textbf{A broad same-input comparison.} Eight systems are evaluated on identical request-side evidence, including repository rules, a character n-gram classifier, untuned zero-shot and few-shot LLM baselines, LoRA and merged artifacts, and f16 and Q4 GGUF paths.
    \item \textbf{Deployment-relevant negative results.} We quantify high safe FPR, complete abstention failure, cross-domain degradation, quantization trade-offs, and message-serialization sensitivity. These findings define where the model is useful and where it is unsafe.
\end{enumerate}

The remainder of the paper reviews related work, defines the task and model, documents data and evaluation controls, presents comparative and ablation results, and discusses the implications for shadow-mode and fail-closed system design.
